\section*{Bab \ref{ch:b1:e-ph-diagrams} --- Diagram E--pH}

\begin{solution}{exo:b1:e-ph-diagrams:1}
Mangan: \ce{Mn} (0) $<$ \ce{Mn^2+}, \ce{Mn(OH)2} ($+$II) $<$ \ce{MnO2}
($+$IV) $<$ \ce{MnO4-} ($+$VII). Klor: \ce{Cl-} ($-$I) $<$ \ce{Cl2} (0)
$<$ \ce{HClO}, \ce{ClO-} ($+$I) $<$ \ce{ClO3-} ($+$V).
\end{solution}

\begin{solution}{exo:b1:e-ph-diagrams:2}
$-0.059 \times 3/1 = -0.177$; $+0.059 \times 2/2 = +0.059$;
$-0.059 \times 8/5 = -0.094$; $-0.059 \times 2/2 = -0.059$ V per satuan
pH. Untuk pasangan tembaga, proton-protonnya dilepaskan di sisi bentuk
tereduksi ($m = -2$): potensialnya naik seiring naiknya pH.
\end{solution}

\begin{solution}{exo:b1:e-ph-diagrams:3}
\ce{H+}/\ce{H2} melewati 0~V pada pH 0 dan tidak pernah mencapai 0.50~V
untuk $\mathrm{pH} \geqslant 0$. \ce{O2}/\ce{H2O} mencapai 0.50~V pada
$\mathrm{pH} = 0.73/0.059 = 12.4$ dan 0~V hanya pada pH 20.8, di luar
skala. Pita itu selebar \qty{1.23}{V} pada setiap pH, termasuk pH 7
(dari $-0.41$ sampai $+0.82$~V).
\end{solution}

\begin{solution}{exo:b1:e-ph-diagrams:4}
(1, 1.0~V): \ce{Fe^3+}. (4, 0): garis \ce{Fe(OH)3}/\ce{Fe^2+} berada pada
$1.09 - 0.71 = 0.38$~V dan titik itu terletak di bawahnya, di atas
$-0.47$~V: \ce{Fe^2+}. (10, $-0.4$~V): antara $-0.07 - 0.59 = -0.66$ dan
$0.28 -
0.59 = -0.31$~V: \ce{Fe(OH)2}. (10, 0.5~V): \ce{Fe(OH)3}.
\end{solution}

\begin{solution}{exo:b1:e-ph-diagrams:5}
\ce{Fe^2+}/\ce{Fe}: $-0.41 + 0.030 \times (-6) = -0.59$~V.
\ce{Fe^3+}/\ce{Fe(OH)3}: $14.00 - \frac13(38.55 - 6) = 3.15$. Daerah-daerah
ion terlarut (korosi) membesar: \ce{Fe^2+} meluas sampai $-0.59$~V dan
\ce{Fe^3+} sampai pH 3.15; logam dan hidroksida kehilangan wilayah.
\end{solution}

\begin{solution}{exo:b1:e-ph-diagrams:6}
\ce{Fe(OH)3 + 3H+ + e- -> Fe^2+ + 3H2O}: $E = E^\circ + 0.059\log\frac{[\ce{H+}]^3}{[\ce{Fe^2+}]}
= (E^\circ - 0.059\log c) - 0.177\,\mathrm{pH}$. Pada pH 1.82 garis itu
bertemu \ce{Fe^3+}/\ce{Fe^2+} pada 0.77~V, sehingga tetapannya
$0.77 + 0.177 \times 1.82
= 1.09$~V.
\end{solution}

\begin{solution}{exo:b1:e-ph-diagrams:7}
$E^\circ(\ce{Cu+}/\ce{Cu}) = 0.52 > E^\circ(\ce{Cu^2+}/\ce{Cu+}) = 0.16$:
menurut \cref{prop:b1:e-ph-diagrams:disproportionation-criterion} \ce{Cu+}
mengalami \omterm{def:b1:e-ph-diagrams:disproportionation}{disproporsionasi}. Satu-satunya batas: $E = 0.34 + 0.030\log 0.010 =
0.28$~V.
\end{solution}

\begin{solution}{exo:b1:e-ph-diagrams:8}
Pada pH 0 daerah tembaga (di bawah 0.28~V) bertumpang tindih dengan
daerah air dan terletak di atas garis hidrogen: tidak ada reaksi dengan
asam klorida tanpa udara. Asam nitrat: \ce{NO3-}/\ce{NO} pada 0.96~V
terletak di atas daerah tembaga, yang teroksidasi:
\ce{3Cu + 2NO3- + 8H+ -> 3Cu^2+ + 2NO + 4H2O}.
% ledger: dfg:NO3-_ao, dfg:NO_g (E0 computed in figdata/bachelor-1/standard-potentials.py)
\end{solution}

\begin{solution}{exo:b1:e-ph-diagrams:9}
Pada pH 3 garis \ce{O2}/\ce{H2O} berada pada $1.23 - 0.18 = 1.05$~V, di
atas seluruh daerah \ce{Fe^2+}: oksigen mengoksidasi besi(II). Di atas
pH 1.82 besi(III) yang terbentuk mengendap sebagai hidroksida yang kuning
kecokelatan: \ce{4Fe^2+ + O2 + 10H2O -> 4Fe(OH)3 + 8H+}.
\end{solution}

\begin{solution}{exo:b1:e-ph-diagrams:10}
Garis-garis tegak: $14.00 - \frac12(16.38 - 2) = 6.81$; dan, dari
\ce{Zn(OH)2 + 2OH- <=> Zn(OH)4^2-} ($K = 10^{-1.93}$) dengan
$[\ce{Zn(OH)4^2-}] = 10^{-2}$, $[\ce{OH-}]^2 = 10^{-0.07}$, pH 13.97.
\ce{Zn^2+}/\ce{Zn}: $-0.76 + 0.030 \times (-2) = -0.82$~V. \ce{Zn(OH)2}/\ce{Zn}:
kemiringan $-0.059$, melalui $(6.81, -0.82)$: $E = -0.42 - 0.059\,\mathrm{pH}$.
\ce{Zn(OH)4^2- + 4H+ + 2e- -> Zn + 4H2O}: kemiringan $-0.118$, melalui
$(13.97, -1.24)$: $E = 0.41 - 0.118\,\mathrm{pH}$.
\end{solution}

\begin{solution}{exo:b1:e-ph-diagrams:11}
Daerah besi terletak di bawah garis hidrogen pada setiap pH. pH 2: besi
teroksidasi menjadi \ce{Fe^2+} dengan pelepasan hidrogen (korosi). pH 7:
reaksi itu masih mungkin, tetapi garis hidrogen ($-0.41$~V) dekat dengan
\ce{Fe^2+}/\ce{Fe} dan reaksinya sangat lambat. pH 13: besi teroksidasi
menjadi \ce{Fe(OH)2}, sebuah \omterm{def:b1:e-ph-diagrams:corrosion-domains}{daerah pasivasi}: terbentuk lapisan dan paku
itu tetap mengilap. Dengan oksigen terlarut, garis \ce{O2}/\ce{H2O}, yang
jauh di atasnya, mengoksidasi besi menjadi besi(III) pada setiap pH:
karat, \ce{Fe(OH)3} dan oksida-oksida dehidratnya.
\end{solution}

\begin{solution}{exo:b1:e-ph-diagrams:12}
Seng mempunyai daerah yang lebih rendah: jika bersentuhan dengan baja,
seng teroksidasi lebih dahulu, \ce{2Zn + O2 + 2H2O -> 2Zn(OH)2} pada
pH 8, dan elektron-elektron yang dilepaskannya mereduksi oksigen di
permukaan baja, yang tetap berupa logam. Magnesium ($E^\circ = -2.36$~V)
lebih rendah lagi dan juga berhasil; tembaga (0.34~V) berada di atas
besi: jika dihubungkan dengan tembaga, besilah logam yang teroksidasi,
lebih cepat daripada jika sendirian.
\end{solution}

\begin{solution}{pb:b1:e-ph-diagrams:1}
\textbf{1.} $-$I, 0, $+$I, $+$I.
\textbf{2.} \ce{Cl-} di bawah, \ce{Cl2}, lalu \ce{HClO} dan \ce{ClO-} di
atas.
\textbf{3.} \ce{HClO} dan \ce{ClO-}, sebuah \omterm{def:b1:predominant-reaction:bronsted}{pasangan asam--basa}.
\textbf{4.} \ce{Cl2 + 2e- -> 2Cl-}.
\textbf{5.} \ce{2HClO + 2H+ + 2e- -> Cl2 + 2H2O}.
\textbf{6.} \ce{ClO- + H2O + 2e- -> Cl- + 2OH-}.
\textbf{7.} \ce{HClO} di bawah pH 7.55, \ce{ClO-} di atasnya.
\textbf{8.} Tidak ada elektron yang dipertukarkan; batasnya ditetapkan
oleh $K_a$ saja.
\textbf{9.} $1/(1 + 10^{7.4 - 7.55}) = 0.59$: \qty{59}{\percent}.
\textbf{10.} $1/(1 + 10^{0.45}) = 0.26$: tiga perempat klor aktif akan
berupa \ce{ClO-}, disinfektan yang lebih lemah.
\textbf{11.} \ce{ClO-}.
\textbf{12.} $E = 1.396 + 0.0296\log\frac{[\ce{Cl2}]}{[\ce{Cl-}]^2} = 1.396
+ 0.0296\log\frac{0.005}{10^{-4}} = 1.396 + 0.050 = 1.446$~V; tidak ada
\ce{H+} dalam \omterm{def:b1:nernst:couple}{setengah reaksinya}.
\textbf{13.} $E = 1.594 + 0.0296\log\frac{[\ce{HClO}]^2[\ce{H+}]^2}{[\ce{Cl2}]}
= 1.594 + 0.0296\log\frac{10^{-4}}{0.005} - 0.0592\,\mathrm{pH} = 1.594 -
0.050 - 0.0592\,\mathrm{pH}$.
\textbf{14.} Dengan menyamakannya, $0.0592\,\mathrm{pH} = 1.594 - 1.396 - 2 \times 0.0503
= 0.0974$, $\mathrm{pH} = 1.646$.
\textbf{15.} \ce{HClO + H+ + 2e- -> Cl- + H2O}: satu proton untuk dua
elektron, kemiringan $-0.0592/2 = -0.0296$.
\textbf{16.} Dengan menjumlahkan kedua \omterm{def:b1:nernst:couple}{setengah reaksi} (masing-masing
satu elektron per atom klor): $2E^\circ = 1.594 + 1.396$,
$E^\circ(\ce{HClO}/\ce{Cl-}) =
1.495$~V; batasnya $E = 1.495 - 0.0296\,\mathrm{pH}$.
\textbf{17.} \ce{ClO- + 2H+ + 2e- -> Cl- + H2O}: $E = 1.718 - 0.0592\,\mathrm{pH}$
(tetapan dipilih demi kekontinuan: $1.495 - 0.0296 \times 7.55 = 1.272 = 1.718
- 0.0592 \times 7.55$).
\textbf{18.} \ce{Cl-} di bawah garis-garis itu; \ce{Cl2} dalam sebuah
segitiga kecil pada $\mathrm{pH} < 1.65$ di atas 1.446~V; \ce{HClO} di
atas garis berkemiringan $-0.0296$ sampai pH 7.55, \ce{ClO-} sesudahnya.
Garis \ce{O2}/\ce{H2O}, dari 1.23 sampai 0.40~V, terletak di bawah semua
batas ini (diagram yang dihitung oleh skrip bab ini mempunyai bentuk
yang sama).
\textbf{19.} Tidak: daerahnya terletak di atas garis oksigen, sehingga
asam itu dapat mengoksidasi air, tetapi reaksinya lambat. Sebuah kolam
kehilangan klornya terutama karena klor itu mengoksidasi apa yang
dibawa perenang dan udara, serta karena sinar matahari; klornya harus
ditambah terus.
\textbf{20.} \ce{Cl2} tidak mempunyai daerah di sana: klor mengalami
\omterm{def:b1:e-ph-diagrams:disproportionation}{disproporsionasi}, \ce{Cl2 + H2O -> HClO + Cl- + H+}; dalam basa
\ce{Cl2 + 2OH- -> ClO- + Cl- + H2O}.
\textbf{21.} Dengan mengalirkan klor ke dalam larutan natrium hidroksida
(reaksi pada pertanyaan 20).
\textbf{22.} \ce{HClO + Cl- + H+ -> Cl2 + H2O}, sebuah \omterm{def:b1:e-ph-diagrams:disproportionation}{komproporsionasi}.
\textbf{23.} Di bawah pH 1.6, hipoklorit dan klorida mengalami
\omterm{def:b1:e-ph-diagrams:disproportionation}{komproporsionasi}: klor terbentuk, dan melewati \omterm{def:b1:precipitation:solubility}{kelarutannya} klor itu
lepas sebagai gas beracun.
\textbf{24.} Pada $c = \qty{0.010}{mol/L}$, pH 4 berada di luar daerah
\ce{Cl2}. Namun, batas itu bergantung pada $c$: mengulangi pertanyaan
12--14 untuk $c$ umum menghasilkan $\mathrm{pH} = 3.34 + \log(2c)$, yang
mencapai 3.6 untuk $c = \qty{1}{mol/L}$, orde konsentrasi pemutih.
Mencampur produk-produk pekat di dekat pH 4 juga tidak aman.
\textbf{25.} \textbf{pH 1.6}: di atasnya, pada $c = \qty{0.010}{mol/L}$,
klor terlarut mengalami \omterm{def:b1:e-ph-diagrams:disproportionation}{disproporsionasi} menjadi asam hipoklorit dan
klorida.
\end{solution}
